\documentclass{situationist-preprint}

\title{Against Infinitism}
\msid{EPI-05-b}
\msfamily{EPI-05}
\msdiscipline{Epistemology}
\mstype{Companion manuscript}
\mskeywords{infinitism; regress of reasons; arbitrariness; experience; structure of justification}

\begin{document}
\maketitle

\begin{abstract}
Infinitism holds that a belief is justified only if there is available to the subject an infinite, non-repeating chain of reasons for it. Its chief argument is that any other view makes the end of the chain arbitrary. We argue, using the tree structure defended in EPI-05-a, that the arbitrariness objection does not touch a view on which the chain ends in experiences, which are not reasons of the kind that could call for further reasons; and that infinitism itself cannot distinguish justified from unjustified beliefs, since chains of available reasons exist for almost any proposition and are not justified by their length.
\end{abstract}

\section{The arbitrariness objection}

\citet{Klein1999} argued that a belief is justified only if the subject has available an infinite and non-repeating series of reasons for it. Against foundationalism he pressed the objection from arbitrariness: if a chain of reasons ends in a basic belief, the subject has no reason for that belief, and holding it is arbitrary; a foundationalist who supplies a reason, such as that beliefs of this kind are likely to be true \citep{BonJour1985}, has extended the chain.

EPI-05-a holds that justification has the structure of a tree of finite height whose leaves are beliefs grounded in experiences (its Theorem~4.3). The objection does not apply to it. What the objection targets is a belief held for no reason. On the tree structure no belief is held for no reason: every leaf belief is held on the basis of an experience that makes its content seem true (Definition~3.1), and is justified only if undefeated (Assumption~3.4). The experience is not a further belief, and so not a further link of the kind Klein's regress demands; it is neither justified nor unjustified, and there is nothing it could be arbitrary to hold. The question \enquote{why believe that the wall is red?} has the answer \enquote{because it looks red}; the question \enquote{why does it look red?} is not a request for justification (EPI-05-a, Section~4.1).

\section{The problem for infinitism}

\begin{proposition}\label{p:idle}
If a chain of beliefs none of which is grounded in experience can justify by its length, then justification does not distinguish between incompatible propositions supported by chains of equal length.
\end{proposition}

\begin{proof}
For any proposition $q$ and infinite sequence $r_1, r_2, \ldots$ of distinct propositions, the chain $q \leftarrow (q \wedge r_1) \leftarrow (q \wedge r_1 \wedge r_2) \leftarrow \cdots$ is infinite and non-repeating, and each member entails its predecessor. Such a chain exists for $q$ and for $\neg q$ alike.
\end{proof}

Klein's reply is that reasons must be available to the subject in a sense that requires more than existence. Every substantive account of availability, however, adds a condition that some links be supported by something other than further links, such as their seeming true to the subject or their being reliably formed. A chain that satisfies such a condition owes its justificatory force to that condition. By Corollary~4.6 of EPI-05-a, an infinite chain none of whose members is grounded in experience justifies none of them; the justification that infinitism ascribes to the chain comes, where it exists, from the experiential grounding of some of its members.

\begin{objection}
Finite minds cannot hold infinitely many reasons, but infinitism requires only that they be available, and the regress the tree structure stops is stopped by fiat.
\end{objection}

\begin{reply}
The regress stops where support passes from belief to experience, which is fixed by what the subject's beliefs are based on, not by a decision of the theorist. That a subject could go on producing reasons does not make the further reasons part of what justifies her belief now; on the tree structure they are new beliefs, with support trees of their own.
\end{reply}

\begin{thebibliography}{9}\small
\bibitem[BonJour(1985)]{BonJour1985} BonJour, L. (1985). \emph{The Structure of Empirical Knowledge}. Harvard University Press.
\bibitem[Klein(1999)]{Klein1999} Klein, P.~D. (1999). Human knowledge and the infinite regress of reasons. \emph{Philosophical Perspectives}, 13, 297--325.
\end{thebibliography}

\end{document}
